% q0002.tex — TEMPLATE dissertativa
% Para questões dissertativas: deixe choices={} vazio.

\question{
  id        = {0002},
  title     = {Fiscal Multiplier},
  source    = {IEO},
  year      = {2022},
  round     = {Final},
  language  = {English},
  topic     = {Macro},
  subtopic  = {Fiscal Policy},
  type      = {dissertative},
  statement = {%
    Suppose a closed economy has a marginal propensity to consume (MPC) of 0.8
    and the government increases its spending by \$500 million.

    \begin{enumerate}[label=(\alph*)]
      \item Calculate the fiscal multiplier and the total change in GDP.
      \item Explain how the presence of income taxes would affect your answer.
    \end{enumerate}
  },
  choices   = {},   % leave empty for dissertative questions
  answer    = {},   % leave empty or write a brief key if applicable
  solution  = {%
    \textbf{(a)} The fiscal multiplier in a closed economy without taxes is:
    \[
      k = \frac{1}{1 - \text{MPC}} = \frac{1}{1 - 0.8} = 5
    \]
    Total change in GDP:
    \[
      \Delta Y = k \cdot \Delta G = 5 \times \$500\text{M} = \$2{,}500\text{ million}
    \]

    \textbf{(b)} With proportional income taxes at rate $t$, the multiplier becomes:
    \[
      k = \frac{1}{1 - \text{MPC}(1-t)}
    \]
    Since $\text{MPC}(1-t) < \text{MPC}$, the denominator increases and the
    multiplier falls. Taxes act as an \emph{automatic stabilizer}, dampening
    both expansions and contractions.
  },
}
